% V. Методика за експериментална проверка.
% Този раздел описва ПЛАНА на експеримента. Числа тук няма и не трябва да се появяват,
% преди да са измерени. Всяко незапълнено място е \TODO.
\section{Методика за експериментална проверка}
\label{sec:method}

\subsection{Измервани величини}

Измерват се две величини. Първата е времето за сходимост, дефинирано като интервалът от
момента на отказа на връзката до момента, в който всички възли, които остават свързани,
притежават маршрутни таблици, съответстващи на новата топология. Втората е наличието и
продължителността на преходни цикли, тоест интервали, в които следващите стъпки на два
или повече възела образуват затворен път. Двете величини се извличат от потока от
събития след изпълнението и не изискват достъп до вътрешното състояние по време на
изпълнение.

\subsection{Топологии и сценарии}

Експериментът се провежда върху набор от топологии с нарастващ диаметър и нарастваща
свързаност, защото и двете величини се очаква да зависят точно от тези два параметъра.
Сценарият е един и същ за всяка топология: мрежата достига установено състояние, отказва
се една връзка, изпълнението продължава до ново установено състояние.
\TODO{да се фиксира наборът от топологии: брой възли, диаметър, степен на свързаност}
\TODO{да се фиксира кои връзки отказват във всяка топология и на какво основание са
избрани}

\subsection{Управлявани параметри}

Управляваните параметри са интервалът за поддържане на съседството, интервалът за
обявяване на отпадане, закъснението на връзките и началното зърно на генератора на
псевдослучайни числа. Всяко изпълнение записва пълния набор от параметри заедно с
резултата, за да бъде повторимо.
\TODO{да се фиксират стойностите и обхватите на управляваните параметри}

\subsection{Условия за валидност}

Резултат се приема за валиден само ако са изпълнени следните условия. Изпълнение с едни
и същи параметри и едно и също зърно дава идентичен поток от събития. Мрежата достига
установено състояние преди отказа и след него. Крайните маршрутни таблици съвпадат с
изчислените независимо върху пълния граф. Изпълнение, което не удовлетворява някое от
условията, се отчита като невалидно и не влиза в резултатите.

\subsection{Форма на резултатите}

Резултатите се представят в две таблици и една фигура. Първата таблица дава времето за
сходимост по топология и по стойност на интервала за отпадане. Втората дава броя и
продължителността на преходните цикли за същите комбинации. Фигурата представя времето
за сходимост като функция от диаметъра на топологията, с легенда по стойност на
интервала за отпадане.
\TODO{броят на повторенията за всяка комбинация се фиксира след пробно изпълнение}
